\DocumentMetadata{
  pdfversion=2.0,
  pdfstandard=ua-2,
  lang=en-US,
  tagging=on,
  tagging-setup={math/setup=mathml-SE,tabsorder=structure}
}
\documentclass[11pt,letterpaper]{article}
\usepackage[margin=0.68in,includefoot]{geometry}
\usepackage{newcomputermodern}
\usepackage{amsmath,amscd,mathtools}
\usepackage{tikz,adjustbox}
\providecommand{\square}{\mdlgwhtsquare}
\usepackage[hidelinks]{hyperref}
\hypersetup{
  pdftitle={Analysis PhD Exam, May 2009},
  pdfauthor={University of Florida Department of Mathematics},
  pdfsubject={Graduate mathematics examination},
  pdfdisplaydoctitle=true,
  bookmarksopen=true
}
\setcounter{secnumdepth}{0}
\setlength{\parindent}{0pt}
\setlength{\parskip}{0.28em}
\setlength{\emergencystretch}{3em}
\setlength{\leftmargini}{2.15em}
\setlength{\leftmarginii}{2.65em}
\setlength{\leftmarginiii}{3em}
\pagestyle{empty}
\raggedbottom
\allowdisplaybreaks
\NewTaggingSocketPlug{block/list/label}{examproblem}{%
  \tagstructbegin{tag=Lbl}\tagstructend%
  \tagstructbegin{tag=\LBody}%
  \tagstructbegin{tag=\ProblemHeadingTag,title-o={\ProblemHeadingText}}%
  \tagmcbegin{tag=\ProblemHeadingTag,actualtext={\ProblemHeadingText}}%
  #2%
  \tagmcend%
  \tagstructend%
}
\newenvironment{examproblems}[2]{%
  \def\ProblemHeadingTag{#2}%
  \AssignTaggingSocketPlug{block/list/label}{examproblem}%
  \begin{enumerate}%
  \setcounter{enumi}{#1}%
  \renewcommand{\labelenumi}{\textbf{\theenumi.}}%
  \setlength{\topsep}{0.35em}%
  \setlength{\partopsep}{0pt}%
  \setlength{\itemsep}{0.48em}%
  \setlength{\parsep}{0.18em}%
}{\end{enumerate}}
\newenvironment{examsubparts}{%
  \AssignTaggingSocketPlug{block/list/label}{default}%
  \begin{enumerate}%
  \setlength{\topsep}{0.18em}%
  \setlength{\partopsep}{0pt}%
  \setlength{\itemsep}{0.16em}%
  \setlength{\parsep}{0.1em}%
}{\end{enumerate}}
\newenvironment{examinstructions}{%
  \par\begingroup\itshape
}{%
  \par\endgroup\vspace{0.18em}
}
\makeatletter
\renewcommand\subsection{\@startsection{subsection}{2}{0pt}%
  {-1.25ex plus -.25ex minus -.1ex}%
  {0.55ex plus .1ex}%
  {\normalfont\large\bfseries}}
\renewcommand{\@maketitle}{%
  \newpage\null\vskip -1em%
  {\footnotesize\bfseries
    \href{https://www.ufl.edu/}{University of Florida}, \href{https://math.ufl.edu/}{Department of Mathematics}\par}%
  \vskip -0.5em%
  \tagstructbegin{tag=H1,title={Analysis PhD Exam, May 2009}}%
  \tagpdfparaOff%
  \tagmcbegin{tag=H1}%
    {\LARGE\bfseries\href{https://gma.math.ufl.edu/exams/analysis-phd/}{Analysis PhD Exam}, May 2009\par}%
  \tagmcend%
  \tagpdfparaOn%
  \tagstructend%
  \vskip 0.25em%
  \tagmcbegin{artifact}%
    \hrule height 1pt%
  \tagmcend%
  \vskip 1em%
}
\makeatother
\title{Analysis PhD Exam, May 2009}
\author{}
\date{}
\begin{document}
\maketitle
\thispagestyle{empty}
\begin{examproblems}{0}{H2}
\def\ProblemHeadingText{Problem 1.}
\item
State and prove the Fubini theorems.
\def\ProblemHeadingText{Problem 2.}
\item
This problem has three parts.
\begin{examsubparts}
\item[\textnormal{(a)}]
State the Vitali integral convergence theorem.
\item[\textnormal{(b)}]
State the Lebesgue Dominated Convergence Theorem.
\item[\textnormal{(c)}]
Use (a) to prove (b) [the convergence in measure version].
\end{examsubparts}
\def\ProblemHeadingText{Problem 3.}
\item
This problem has two parts.
\begin{examsubparts}
\item[\textnormal{(a)}]
Let \(\mu:\Sigma\to[-\infty,\infty]\) be countably additive, where~\(\Sigma\) is a~\(\sigma\)-algebra. Define positive and negative sets for~\(\mu\) and state the Hahn decomposition theorem for~\(\mu\).
\item[\textnormal{(b)}]
Let~\(\mu\) be as above. Does there exist a set \(E\in\Sigma\) such that \(\mu(E)\) is the minimum value for~\(\mu\) on~\(\Sigma\)? Prove.
\end{examsubparts}
\def\ProblemHeadingText{Problem 4.}
\item
Suppose \(\{f_n\}\) is a sequence of integrable functions with the property that \(\sum \int |f_n|\,d\mu<\infty\). Analyze the pointwise convergence of \(\sum f_n(x)\), for \(x\in X\).
\def\ProblemHeadingText{Problem 5.}
\item
Let~\(Y\) be an orthonormal set in a Hilbert space and suppose \(f:Y\to\mathbb C\) is given. State necessary and sufficient conditions for \(\sum_{y\in Y}f(y)y\) to converge. Prove.
\def\ProblemHeadingText{Problem 6.}
\item
Let~\(X\) be a normed space and \(\pi:X\to X^{**}\) the natural map. Under what conditions is \(\pi(X)\) closed in \(X^{**}\) (equipped with the norm topology)? Prove.
\def\ProblemHeadingText{Problem 7.}
\item
This problem has two parts.
\begin{examsubparts}
\item[\textnormal{(a)}]
State the Principle of Uniform Boundedness for a family of operators.
\item[\textnormal{(b)}]
Let~\(X\) be a normed space. Use (a) to show that if \(A\subset X\) satisfies \(\sup_{a\in A}|x^*a|<\infty\) for each \(x^*\in X^*\), then \(\sup_{a\in A}|a|<\infty\).
\end{examsubparts}
\def\ProblemHeadingText{Problem 8.}
\item
Let~\(X\) be a normed space, and assume~\(M\) is a closed linear subspace of~\(X\). Let \(z\in X\setminus M\) and \(S=\operatorname{span}\{z,M\}\). Show~\(S\) is closed in~\(X\). [Hint: Let \(f:S\to\mathbb C\) be defined by \(f(x+\alpha z)=\alpha\), for \(x\in M\) and \(\alpha\in\mathbb C\). Show \(\|f\|\le 1/\operatorname{dist}(z,M)\), hence \(f\in S^*\).]
\end{examproblems}
\end{document}
